\section*{Chapter \ref{ch:g7:earthquakes} --- Earthquakes}

\begin{solution}{exo:g7:earthquakes:1}
(a) \omterm{def:g7:earthquakes:focus}{focus}; (b) \omterm{def:g7:earthquakes:intensity}{intensity}; (c) \omterm{def:g7:earthquakes:fault}{fault}; (d) \omterm{def:g7:earthquakes:focus}{epicentre}; (e) \omterm{def:g7:earthquakes:magnitude}{magnitude}; (f)
\omterm{def:g7:earthquakes:seismometer}{seismometer}.
\end{solution}

\begin{solution}{exo:g7:earthquakes:2}
(a) At the village itself: the \omterm{def:g7:earthquakes:focus}{epicentre} is the point of the surface
straight above the \omterm{def:g7:earthquakes:focus}{focus}. (b) \qty{30}{km}, straight down. (c) Usually
lower: the shaking weakens with distance. (It also depends on the ground:
on soft \omterm{def:g3:stones-pebbles-sand:sand}{mud}, the second village could still shake hard.)
\end{solution}

\begin{solution}{exo:g7:earthquakes:3}
(a) \omterm{def:g7:earthquakes:magnitude}{Magnitude} (a number with one decimal). (b) \omterm{def:g7:earthquakes:intensity}{Intensity} (what was seen
in one place). (c) \omterm{def:g7:earthquakes:intensity}{Intensity} (Roman numerals, different in two places).
(d) \omterm{def:g7:earthquakes:magnitude}{Magnitude}: the energy released, one step of \omterm{def:g7:earthquakes:magnitude}{magnitude} more.
\end{solution}

\begin{solution}{exo:g7:earthquakes:4}
(a) False: scientists cannot predict \omterm{def:g4:volcanoes-earthquakes-first-look:earthquake}{earthquakes}; they only know which
\omterm{def:g7:earthquakes:fault}{faults} will break again one day. (b) False: in the open \omterm{def:g2:water-on-earth:ocean}{ocean} a \omterm{def:g7:earthquakes:tsunami}{tsunami}
is rarely more than a metre high; it grows near the \omterm{def:g1:landscapes-around-us:water}{coast}. (c) True. (d)
True: 92~\% of the \omterm{def:g4:volcanoes-earthquakes-first-look:earthquake}{earthquakes} of \omterm{def:g7:earthquakes:magnitude}{magnitude}~6 or more happen less than
\qty{300}{km} from a \omterm{def:g7:plate-tectonics:plate-boundary}{plate boundary}.
\end{solution}

\begin{solution}{exo:g7:earthquakes:5}
(a) About 127 a year of \omterm{def:g7:earthquakes:magnitude}{magnitude} 6 to 6.9; about 13 of \omterm{def:g7:earthquakes:magnitude}{magnitude} 7 to
7.9. (b) $126.7 \div 13.12 \approx 9.7$: about 10~times fewer. (c)
$13.12 + 0.76 + 0.04 = 13.92$: about \textbf{14} a year. (d)
$0.04 \times 50 = 2$: two in 50~years (Sumatra 2004 and Japan 2011).
\end{solution}

\begin{solution}{exo:g7:earthquakes:6}
(a) One step: 32~times. (b) Two steps: $32 \times 32 = 1024$, about
1000~times. (c) Each step makes the waves 10~times larger: $10 \times 10 =
100$~times larger.
\end{solution}

\begin{solution}{exo:g7:earthquakes:7}
(a) $8 \times 5 = \qty{40}{km}$. (b) $8 \times 25 = \qty{200}{km}$. (c)
$120 \div 8 = \qty{15}{s}$. (d) One station gives only a distance, not a
direction: all the points at that distance form a circle around the
station, and the \omterm{def:g7:earthquakes:focus}{epicentre} could be anywhere on it.
\end{solution}

\begin{solution}{exo:g7:earthquakes:8}
(a) V. (b) II. (c) VIII. (d) The \omterm{def:g7:earthquakes:intensity}{intensity} changes from place to place:
report (c) comes from a place close to the \omterm{def:g7:earthquakes:focus}{epicentre}, or built on soft
ground; report (a) from a place farther away, or on solid \omterm{def:g3:stones-pebbles-sand:rock}{rock}. The
\omterm{def:g4:volcanoes-earthquakes-first-look:earthquake}{earthquake} still has only one \omterm{def:g7:earthquakes:magnitude}{magnitude}.
\end{solution}

\begin{solution}{exo:g7:earthquakes:9}
(a) $7031 - 5561 = 1470$. (b) $1470 \div 7031 \approx 0.209$: about
\textbf{21~\%}. (c) The sinking \omterm{def:g7:plate-tectonics:plate}{plate} goes down deep into the \omterm{def:g6:inside-the-earth:layers}{mantle}. It
is still rigid enough to break, so \omterm{def:g4:volcanoes-earthquakes-first-look:earthquake}{earthquakes} can start inside it,
hundreds of kilometres down.
\end{solution}

\begin{solution}{exo:g7:earthquakes:10}
(a) $2500 \div 800 = 3.125$~hours: 3~hours and $0.125 \times 60 = 7.5$
minutes, about 3~hours 8~minutes. (b) 3~h~8~min $-$ 2~h~17~min $=$
51~minutes later. (c) The wave took hours to arrive, while a large
\omterm{def:g4:volcanoes-earthquakes-first-look:earthquake}{earthquake} under the \omterm{def:g2:water-on-earth:ocean}{sea} is located within minutes: a warning could have
sent people to high ground long before the wave came.
\end{solution}

\begin{solution}{exo:g7:earthquakes:11}
(a) $400 \div 800 = 0.5$~hour: 30~minutes. (b) $20 \div 40 = 0.5$~hour:
30~minutes. (c) 60~minutes. (d) The last \qty{20}{km} take as long as the
first \qty{400}{km}: the wave slows down enormously in shallow water. (e)
The front of the wave slows down while the water behind it still comes
fast: the water piles up, and the wave grows taller.
\end{solution}

\begin{solution}{exo:g7:earthquakes:12}
(a) Three steps: $32 \times 32 \times 32 = 32\,768$, about
\textbf{33\,000}~times. (b) $32\,768 \div 127 \approx 258$: about
260~years of all the magnitude-6 \omterm{def:g4:volcanoes-earthquakes-first-look:earthquake}{earthquakes} of the world. (c) Almost all
the energy is released by the few very largest \omterm{def:g4:volcanoes-earthquakes-first-look:earthquake}{earthquakes}: a single
\omterm{def:g7:earthquakes:magnitude}{magnitude}~9 outweighs centuries of magnitude-6 \omterm{def:g4:volcanoes-earthquakes-first-look:earthquake}{earthquakes}.
\end{solution}

\begin{solution}{pb:g7:earthquakes:1}
\textbf{1.} A: $45 - 25 = 20$~s. B: $29 - 16 = 13$~s. C: $34 - 19 =
15$~s.

\textbf{2.} Station B: its delay is the shortest, and the shorter the
delay, the closer the \omterm{def:g4:volcanoes-earthquakes-first-look:earthquake}{earthquake}. (Its faster waves also arrived first.)

\textbf{3.} A: $8 \times 20 = \qty{160}{km}$. B: $8 \times 13 =
\qty{104}{km}$. C: $8 \times 15 = \qty{120}{km}$.

\textbf{4.} All the points \qty{160}{km} from A: a circle of radius
\qty{160}{km} centred on A. The \omterm{def:g7:earthquakes:focus}{epicentre} is somewhere on it.

\textbf{5.} A: $160 \div 20 = 8$~squares. B: $104 \div 20 = 5.2$~squares.
C: $120 \div 20 = 6$~squares.

\textbf{6.} E is $3 + 6 = 9$~squares east and 14.5~squares north of O, out
at \omterm{def:g2:water-on-earth:ocean}{sea} between the three stations; the circles of A and B pass through it.

\textbf{7.} Under the \omterm{def:g2:water-on-earth:ocean}{sea}. A sudden movement of the \omterm{def:g2:water-on-earth:ocean}{sea} floor can set off a
\omterm{def:g7:earthquakes:tsunami}{tsunami}: the \omterm{def:g1:landscapes-around-us:water}{coasts} and the island have to be warned.

\textbf{8.} One step of \omterm{def:g7:earthquakes:magnitude}{magnitude}: about 32~times more energy.

\textbf{9.} $12 \times 20 = \qty{240}{km}$.

\textbf{10.} $240 \div 800 = 0.3$~hour, and $0.3 \times 60 = 18$~minutes.

\textbf{11.} 5~minutes after the \omterm{def:g4:volcanoes-earthquakes-first-look:earthquake}{earthquake}.

\textbf{12.} $18 - 5 = 13$: the people of Port Hale have about
\textbf{13~minutes} to reach high ground. (The wave slows down in the
shallow water near the island, which gives a little more time, but nobody
should count on it: they leave at once.)
\end{solution}
